\documentclass{article}
\usepackage{amsmath}
\usepackage{graphicx}
\usepackage{hyperref}

\title{Regulation of the Lac Operon}
\author{Howard Chen}

\begin{document}
\maketitle

\begin{abstract}
We review transcriptional control of the \textit{lac} operon in \textit{E.\ coli}, including the repressor--operator interaction and catabolite activation.
\end{abstract}

\section{Structure}
The \textit{lac} operon comprises three structural genes (\textit{lacZ}, \textit{lacY}, \textit{lacA}) under the control of a single promoter and operator.

\section{Repression}
In the absence of lactose, LacI tetramers bind the operator region, blocking RNA polymerase.

\section{Induction}
Allolactose (produced by basal $\beta$-galactosidase activity) binds LacI, triggering a conformational change that releases it from the operator.

\section{Catabolite repression}
When glucose is abundant, cAMP levels drop and CRP cannot bind the activator site — the operon is transcribed weakly even if lactose is present.

\section{Conclusion}
The \textit{lac} operon remains a canonical example of combinatorial gene regulation.

\end{document}
